\section{Introduction}\label{sec:introduction}

The word problem for a group with a specified finite generating set asks
whether a given word in the generators and their inverses represents the
identity. It is \emph{decidable} if an algorithm answers this question for
every word. A group is \emph{simple} here if it is nontrivial and its only
normal subgroups are the identity subgroup and the whole group.

\begin{theorem}\label{thm:main}
For every finitely generated group \(G\), the following are equivalent:
\begin{enumerate}[label=\textup{(\roman*)}]
\item \(G\) has decidable word problem.
\item \(G\) embeds in a finitely presented simple group.
\end{enumerate}
\end{theorem}

An embedding throughout means an injective group homomorphism. In the
constructive implication, \(G\) is supplied with a finite generating set
and a word-problem decider; it need not be finitely presented.
Theorem~\ref{thm:main} resolves the Boone--Higman conjecture positively.

Higman's embedding theorem characterizes the finitely generated
subgroups of finitely presented groups as those admitting a recursively
enumerable set of defining relations \cite{Higman61}; see also
\cite[Theorem~A]{AanderaaCohen1980}.
Decidability of the word problem is stronger: it supplies an algorithm
for both equality and inequality. Boone and Higman characterized this
condition by the existence of embeddings
\[
 G\ \leq\ H\ \leq\ K,
 \qquad H\text{ simple},\quad K\text{ finitely presented}
\]
\cite[Theorem~I]{BH74}. They asked whether the simple overgroup \(H\)
could itself be chosen finitely presented \cite[p.~43]{BH74}.
The two properties in their theorem belong to different groups, and
finite presentation does not in general pass to subgroups.
Thompson subsequently obtained a finitely generated simple overgroup
with decidable word problem \cite{Tho80}; a later proof appears in
\cite[Theorem~7.1]{DarbinyanSteenbock2022}.
This strengthens the conclusion to finite generation of the simple
overgroup, but does not supply its finite presentation.

Positive cases of the conjecture include hyperbolic and contracting
self-similar groups \cite[Theorem~A and Corollary~D]{BBMZhyperbolic},
Baumslag--Solitar groups and finite-rank free-by-cyclic groups
\cite[Theorem~A]{BuxLlosaWu2025}, and automorphism groups of finite-rank
free groups and mapping class groups of orientable punctured surfaces
of finite type \cite[Theorem~A and Corollary~B]{BFFHZ}.
The survey \cite{BBMZsurvey} describes further results and the
development of the conjecture.

Our principal intermediate result concerns permutation actions.
An action on a set \(X\) is \emph{faithful} if only the identity acts
trivially, and \emph{two-transitive} if it is transitive on the ordered
pairs of distinct points of \(X\).

\begin{theorem}\label{thm:envelope}
Every finitely generated group with decidable word problem embeds in a
finitely presented group \(\Gamma\) admitting a faithful two-transitive
action on a countably infinite set. Every point stabilizer of this action
is finitely generated.
\end{theorem}

The action is designed for the twisted Brin--Thompson construction of
Belk and Zaremsky, which embeds an acting group in a simple group
\cite[Section~3 and Theorem~3.4]{BelkZaremsky2022}.
Their original finite-presentation theorem used finitely many orbits on
finite subsets of each size and finitely presented finite-set
stabilizers \cite[Theorem~D]{BelkZaremsky2022}.
Zaremsky subsequently showed that finite presentation of the acting
group, finite generation of point stabilizers, and finitely many orbits
on unordered pairs suffice
\cite[Theorem~A and Corollary~B]{Zaremsky2024Finite}.
Two-transitivity supplies the last condition here. We state this
criterion precisely in Theorem~\ref{thm:zaremsky}; the construction below
provides its faithful action and the embedding of the prescribed group.

\subsection*{The proof}

The action is affine. Translations move the first point of an ordered
pair to zero, so two-transitivity reduces to transitivity of the linear
part on nonzero vectors. A simple algebra \(A\) provides the required
coefficient operations: if \(0\ne m\in A\), then its identity belongs
to \(AmA\), and hence is a finite sum of left and right multiples of
\(m\). Elementary row operations can use this identity to move vectors.
Finite presentation, however, will come from a different algebra.

Over \(k=\F_2\), we construct a finitely presented algebra \(B\), a
simple subalgebra \(A\subseteq B\) with its own identity, and an
injective nonunital map \(\phi:B\to B\) with image in \(A\).
The additive direct limit \(M_\infty=\varinjlim(B,\phi)\) is the
carrier for the action on \(X=M_\infty^6\). Every finite collection
of its coordinates lies in one copy of \(B\); advancing once places
them in \(A\). Thus simplicity becomes available at a later stage
without requiring the finitely presented algebra \(B\) itself to be
simple.

The coefficient ring \(R=B\otimes_k B^{\op}\) acts on \(B\) by
left and right multiplication. We give an explicit finite presentation
of the affine Steinberg group
\[
 \Gamma_0=B^6\rtimes\St_6(R).
\]
The Steinberg group maps onto the elementary matrix group \(\E_6(R)\),
but its kernel is an obstruction to the desired faithful affine action.
We remove that entire kernel by a single application of the coefficient
map induced by \(\phi\). The key algebraic input consists of two
splitting pairs commuting with the coefficient image. Their corner
identities give an exact calculation in the Steinberg presentation;
no centrality hypothesis on the kernel is needed.

The resulting endomorphism of \(\Gamma_0\) defines a finitely
presented mapping torus \(\Gamma\). In its direct-limit description
the Steinberg kernel has disappeared, leaving elementary matrices acting
on \(M_\infty^6\), together with translations and a stage shift.
Simplicity of \(A\) gives the nonzero-vector orbit required for
two-transitivity. Faithfulness uses two further properties of the
algebra construction: a separating vector detects every nonzero
coefficient after one advance, and the proper support
\(\phi(1_B)\ne1_B\) makes the stage shift detectable. The stabilizer
of zero is the quotient by the translation subgroup, so it is finitely
generated. To embed the original \(G\), we first use an HNN extension
to express its elements as controlled commutators; balanced elementary
diagonals then place this copy of \(G\) in the linear part.

It remains to construct the algebras with these simultaneous properties.
We do so through a finitely presented monoid and a decidable monoid that
encode one another. A finite presentation is compiled syntactically from
a program index, before any totality assertion is available. Its partial
evaluator controls the length of every call to the encoded program.
Long codes make the calls needed to normalize a word strictly shorter
than that word. A fixed-point argument and induction on length then
establish totality and compatibility with multiplication.
Finite-tuple selectors isolate individual basis terms in a contracted
monoid algebra, producing \(A\); a separate separator distinguishes
ordered pairs of basis terms, producing the faithful coefficient detector.

The finite-presentation construction follows the machine-encoding
tradition of Post and Murskii \cite{Post1947,Murskii1967}.
Here the bounded evaluator permits the two monoids to be constructed
together. The selectors strengthen the contextual-collapse property in
the semigroup counterpart of the Boone--Higman theorem
\cite[pp.~49--52]{BH74}: isolating each term of a finite linear
combination gives the simplicity needed for the affine action.

Section~\ref{sec:reduction} states the algebraic data and gives the
HNN reduction.
Sections~\ref{sec:affine}--\ref{sec:action} construct the action and prove
the two theorems from those data.
Section~\ref{sec:algebra} constructs the algebras from a precise monoid
theorem, proved by bounded self-reference in Section~\ref{sec:bootstrap}.
Appendix~\ref{sec:compiler} gives the full finite compiler and its
evaluator proof.
